%!TEX root = ../../main.tex
% ============================================================
% 数学\几何作图\五年级.tex
% 本文件由 main.tex 通过 \input 引入，请勿单独编译
% ============================================================


\section{解答：图形的平移与面积转化}

\vspace{0.4cm}

\noindent\textbf{4. 把下面平行四边形中的涂色三角形向右平移 5 格，看一看，平移后的三角形与平行四边形中的空白部分组成了一个什么图形？ (\makebox[2cm][c]{\textcolor{blue}{\textbf{长方形}}})}

\vspace{0.8cm}

\begin{center}
\begin{tikzpicture}[>=Stealth, line width=1pt, scale=0.75] % 整体思路：先画方格背景，再画原平行四边形及左侧阴影三角形，接着把这个三角形按“向右平移 5 格”得到新位置，用箭头说明移动过程

  % 外层实线网格边框
  \draw[gray, thin] (0,0) rectangle (11,5); % 画最外层矩形边框，范围是 11×5 的方格区域
  % 内部虚线网格
  \draw[gray, dashed, thin, step=1] (0,0) grid (11,5); % 用步长 1 的虚线网格填满区域；step=1 表示每格边长为 1 个坐标单位
  % 再绘制一次清晰的实线外框以覆盖重叠的虚线边角
  \draw[gray, thin] (0,0) rectangle (11,5); % 再描一次实线边框，使外框不被虚线削弱

  % 原始平行四边形在方格图中的严谨真实坐标
  \coordinate (A) at (2,1); % 定义原平行四边形左下顶点 A
  \coordinate (B) at (7,1); % 定义原平行四边形右下顶点 B
  \coordinate (C) at (9,4); % 定义右上顶点 C，较底边向右偏 2 格、向上 3 格
  \coordinate (D) at (4,4); % 定义左上顶点 D
  \coordinate (H) at (4,1); % 定义高的垂足 H，也是左侧小三角形在底边上的右端点

  % 原涂色三角形 (灰色)
  \fill[black!30] (A) -- (H) -- (D) -- cycle; % 填充左侧直角三角形 AHD，作为要平移的涂色部分
  
  % 原平行四边形黑色实线轮廓
  \draw (A) -- (B) -- (C) -- (D) -- cycle; % 画原平行四边形外轮廓
  % 内部高线 (作为三角形和右侧空白部分的分割线)
  \draw (D) -- (H); % 画内部高线 DH，把平行四边形分成左侧三角形与右侧空白部分
  
  % 原图中直角标记标注在高的右边
  \pic [draw, line width=0.8pt, angle radius=0.25cm] {right angle = D--H--B}; % 在 H 处画直角标记，说明 DH 垂直于底边 HB

  % --- 平移后的目标位置三角形 ---
  % 根据题意严格按比例：横向坐标向右平移 +5，纵向坐标 +0
  \coordinate (A') at (7,1); % 平移后点 A'，即原 A 向右移动 5 格的位置
  \coordinate (H') at (9,1); % 平移后点 H'
  \coordinate (D') at (9,4); % 平移后点 D'，正好与原顶点 C 重合
  
  % 平移辅助箭头
  \draw[->, blue, line width=1.4pt] (3.0, 2.5) -- (8.0, 2.5); % 画蓝色箭头表示平移方向与距离；箭头从左向右
  \node[blue, above, font=\small, fill=white, inner sep=2pt] at (6, 2.5) {向右平移 5 格}; % 在箭头上方放说明文字；fill=white 给节点白底，避免与网格线重叠
  
  % 平移后的新三角形填色与虚线轮廓 (以淡蓝色突显出它的平移结果痕迹)
  \fill[blue!15, opacity=0.8] (A') -- (H') -- (D') -- cycle; % 在平移后位置填充淡蓝色三角形，表示新位置的图形；opacity=0.8 表示略带透明度
  \draw[blue, dashed, line width=1.2pt] (A') -- (H') -- (D') -- cycle; % 用蓝色虚线描出平移后三角形轮廓

\end{tikzpicture}
\end{center}

\vspace{0.6cm}
{\small\color{gray}
\textbf{解析：}\\
观察方格图可以看出，左侧涂色的三角形是一个底边长为 $2$ 格、高为 $3$ 格的直角三角形。平行四边形被高切分开后，其右边空出的部分由于原先是平行倾斜的所以有一个正好也是底向外延展了 $2$ 格、高为 $3$ 格的缺口形补位。\\
当我们把这块涂色的三角形严格沿着格子\textbf{向右平移 5 格}（代码中的坐标精准体现了由横向长度相加 $5$ 推导而来）后，它的垂直边完全平移到了 $x=9$ 的纵线上，而其斜边则地同平行四边形原本右侧的那条倾斜边（从 $(7,1)$ 到 $(9,4)$）重合了。平移补齐后，与原有的右侧空白部分拼放在一起，形成了左右边都为铅质垂线、上下边都在水平网格线上的新四边形。数一数下方的网格可知，这个新图形的长边为平移落差赋予的 $5$ 个单位、短边（宽）为原来的高 $3$ 个单位，且四个角通过平移互补都被凑成了规整的直角，因此这个合并后的图形是一个\textbf{长方形}。
}


\vspace{0.6cm}

{\small\color{blue!70!black}
\textbf{绘图基本原理与步骤：}\\
\textbf{1. 坐标定位与几何构建}：根据题意中的已知长度和角度，使用 \texttt{coordinate} 定义关键顶点坐标，通过 \texttt{draw} 指令按序连接成完整几何图形。线宽、颜色参数区分原图（黑色）与答案（红色 \texttt{red!70!black}）图层。\\
\textbf{2. 辅助量标注}：利用 \texttt{node} 的方向锚点（\texttt{above}、\texttt{below}、\texttt{left}、\texttt{right}）将角度、边长、面积等数据标注在图形周围，避免与线条或其他标注重叠。若涉及角弧则使用 \texttt{arc} 或 \texttt{pic[angle]} 命令渲染。\\
\textbf{3. 虚实线分层}：题干已知条件用实线绘制，解答新增的辅助线或操作痕迹用 \texttt{dashed} 虚线或彩色加粗线标识，确保读者一眼区分原有信息与作答内容。
}

%% ==============================
%% 同底等高的平行四边形画法
%% ==============================
\newpage



\section{解答：画同底等高的平行四边形}

\vspace{0.4cm}

\noindent\textbf{5. 下面的这组平行线之间已经画了一个平行四边形。在这组平行线之间再画两个平行四边形，并使它们与已知的平行四边形等底等高。}

\vspace{0.8cm}

\begin{center}
\begin{tikzpicture}[>=Stealth, line width=1pt] % 整体思路：先画上下两条平行线，再在同一条底边 AB 上画三个顶点不同、但都落在上方平行线上的平行四边形，以体现“同底等高”

  % 辅助平行线
  \draw[gray, thick] (-0.5, 2.5) -- (12.5, 2.5); % 画上方辅助平行线，高度固定为 y=2.5
  \draw[gray, thick] (-0.5, 0) -- (12.5, 0); % 画下方辅助平行线，作为共用底边所在直线

  % 基准坐标
  \coordinate (A) at (2.5, 0); % 共用底边左端点 A
  \coordinate (B) at (6.5, 0); % 共用底边右端点 B，AB 长为 4
  
  % 原先的灰色平行四边形 (根据原图像素严格推导：高2.5, 底宽4, 底高比1.6:1, 基本水平倾斜距离1.25)
  \coordinate (D) at (3.75, 2.5); % 原灰色平行四边形的左上点 D
  \coordinate (C) at (7.75, 2.5); % 原灰色平行四边形的右上点 C
  
  \fill[black!30] (A) -- (B) -- (C) -- (D) -- cycle; % 先填充原有的灰色平行四边形
  \draw (A) -- (B) -- (C) -- (D) -- cycle; % 再描边画出原平行四边形轮廓

  % 第一个补充的平行四边形 (红色)
  \coordinate (D_red) at (5.5, 2.5); % 第一个新增平行四边形的左上点
  \coordinate (C_red) at (9.5, 2.5); % 第一个新增平行四边形的右上点，与 AB 等长
  \draw[red, thick] (A) -- (B) -- (C_red) -- (D_red) -- cycle; % 用红色粗线画出第一个“同底等高”平行四边形

  % 第二个补充的平行四边形 (绿色)
  \coordinate (D_green) at (7.25, 2.5); % 第二个新增平行四边形的左上点
  \coordinate (C_green) at (11.25, 2.5); % 第二个新增平行四边形的右上点
  \draw[green!60!black, thick] (A) -- (B) -- (C_green) -- (D_green) -- cycle; % 用绿色粗线画出第二个“同底等高”平行四边形

\end{tikzpicture}
\end{center}

\vspace{0.6cm}
{\small\color{gray}
\textbf{解析：}\\
平行四边形的面积公式为 $S = \text{底} \times \text{高}$。如果两个平行四边形\textbf{等底等高}，那么它们的面积一定相等。\\
要画完全符合要求且等面积的平行四边形，只需要满足两个条件：\\
1. \textbf{等底}：与已知图形共用一条底边（即上述代码作图时固定了底部线段的两端点坐标不变）。\\
2. \textbf{等高}：顶点只能落在同一条平行线上（即利用题目给定的那一组固定高度差的平行线）。\\
在作图时，我们固定了下方的底边宽，然后在上方的平行线上取了长度与底边严格相同（即宽为 $4$ 厘米比例）的两段线段作为新的上底（如红线上底和绿线上底），并将它们与共用的下底分别连接即可。红、绿两个平行四边形虽然倾斜得很严重，但它们的底宽、高落差均与原来的灰色图形一致，属于“同底等高”的规范图形。
}


\vspace{0.6cm}

{\small\color{blue!70!black}
\textbf{绘图基本原理与步骤：}\\
\textbf{1. 坐标定位与几何构建}：根据题意中的已知长度和角度，使用 \texttt{coordinate} 定义关键顶点坐标，通过 \texttt{draw} 指令按序连接成完整几何图形。线宽、颜色参数区分原图（黑色）与答案（红色 \texttt{red!70!black}）图层。\\
\textbf{2. 辅助量标注}：利用 \texttt{node} 的方向锚点（\texttt{above}、\texttt{below}、\texttt{left}、\texttt{right}）将角度、边长、面积等数据标注在图形周围，避免与线条或其他标注重叠。若涉及角弧则使用 \texttt{arc} 或 \texttt{pic[angle]} 命令渲染。\\
\textbf{3. 虚实线分层}：题干已知条件用实线绘制，解答新增的辅助线或操作痕迹用 \texttt{dashed} 虚线或彩色加粗线标识，确保读者一眼区分原有信息与作答内容。
}

%% ==============================
%% 行程问题：相遇点位置精确作图
%% ==============================
\newpage



\section{解答：面积增加问题——长方形扩建正方形}

\vspace{0.4cm}

\noindent\textbf{6. 赵大伯家有一个长方形鱼塘，长比宽多 20 米。如果把这个鱼塘扩建成为一个最小的正方形，面积就增加 1600 平方米。原来鱼塘的面积是多少平方米？（先根据题意画出示意图，再解答）}

\vspace{0.8cm}

\begin{center}
\begin{tikzpicture}[>=Stealth, line width=1pt, scale=0.08] % 整体思路：把原鱼塘画成上方长方形，把扩建增加的部分画成下方长方形，两块上下拼起来正好组成最小正方形
  % 原始鱼塘 (80x60)
  \draw[thick] (0, 20) rectangle (80, 80); % 画原鱼塘长方形：左下角 (0,20)，右上角 (80,80)，长 80、宽 60
  % 扩建部分 (80x20)
  \draw[thick] (0, 0) rectangle (80, 20); % 画扩建增加的下方长方形：长 80、宽 20
  
  % 标注
  \node[above=5pt, font=\large] at (40, 80) {长}; % 在上边中点上方标“长”
  \node[right=5pt, font=\large] at (80, 50) {宽}; % 在原鱼塘右侧中部标“宽”
  \node[right=5pt, font=\large] at (80, 10) {20米}; % 在新增部分右侧标出宽差 20 米
  \node[font=\large] at (40, 10) {1600平方米}; % 在新增矩形中间标出增加面积 1600 平方米
\end{tikzpicture}
\end{center}

\vspace{0.6cm}

{\small\color{gray}
\textbf{解析：}\\
1. **分析扩建部分**：长方形鱼塘的长比宽多 20 米 ($L = W + 20$)。要扩建为“最小的正方形”，由于长大于宽，正方形的边长应等于原先的长 ($L$)。\\
2. **确定增加面积的形状**：扩建的部分是一个长为 $L$、宽为 20 米 ($L - W$) 的长方形。\\
3. **计算长和宽**：已知增加的面积为 1600 平方米，根据公式 $S = \text{长} \times \text{宽}$，可得 $L = 1600 \div 20 = 80$（米）。\\
4. **求原面积**：原先的宽 $W = 80 - 20 = 60$（米）。原鱼塘的面积为 $80 \times 60 = 4800$（平方米）。
}

\vspace{0.4cm}
\noindent\textbf{解答：}\\
\indent $1600 \div 20 = 80$（米）\dots\dots\dots 原鱼塘的长\\
\indent $80 - 20 = 60$（米）\dots\dots\dots 原鱼塘的宽\\
\indent $80 \times 60 = 4800$（平方米）

\vspace{0.2cm}
\noindent\textbf{答：}原来鱼塘的面积是 4800 平方米。

\vspace{0.6cm}

{\small\color{blue!70!black}
\textbf{绘图基本原理与步骤：}\\
\textbf{1. 几何模型构建}：将题目文字描述的“长方形扩建为正方形”抽象为两个拼接的矩形，大矩形边长代表原长。\\
\textbf{2. 比例缩放}：基于实际长度（80m和60m）设定绘图比例（本图取 1:0.08），确保图形视觉上的长宽关系符合真实逻辑。\\
\textbf{3. 关键标注对齐}：将“数值量”（1600平方米）置于新增面积块中心，“文字量”（长、宽）置于对应边侧，辅助理解代数推导过程。
}


\newpage




\section{解答：在点子图上画面积相等的图形}

\vspace{0.4cm}

\noindent\textbf{1. 在点子图上画和长方形面积相等的平行四边形、三角形和梯形各一个。}

\vspace{1.2cm}

\begin{center}
% =====================================================================
% 【整体绘制思路：等面积图形的构造】
% 1. 点子图背景：
%    使用 \foreach 循环绘制 21x7 的点阵阵列，模拟题目要求的“点子图”环境。
%    点阵原点为 (0,0)，最大至 (20,6)。
% 2. 图形面积逻辑运算 (目标面积为 8)：
%    - 给定长方形：宽 2，高 4。计算面积 = 2 * 4 = 8。
% 3. 构造等积作图图形 (所有高统一取最直观的 4 单位)：
%    - 平行四边形：面积 = 底 * 高。需要 底 * 4 = 8，得出 底 = 2。
%    - 三角形：面积 = 底 * 高 / 2。需要 底 * 4 / 2 = 8，得出 底 = 4。
%    - 梯形：面积 = (上底 + 下底) * 高 / 2。需要 (上底 + 下底) * 4 / 2 = 8，
%            得出 (上底 + 下底) = 4。可以取 上底=1，下底=3。
% 4. 坐标排布规则：
%    - 每个图形之间保持 1~2 个单位的间隔，均匀分布在点阵图上。
%    - 为了清晰，已知图形用黑色表示，要求画出的图形全部用红色表现。
% =====================================================================
\begin{tikzpicture}[>=Stealth, line width=1.0pt, scale=0.8]

  % =========== 背景点子图 ===========
  % \foreach \x ... \foreach \y ... 批量画出所有的点
  \foreach \x in {0,1,2,3,4,5,6,7,8,9,10,11,12,13,14,15,16,17,18,19,20} {
    \foreach \y in {0,1,2,3,4,5,6} {
      \fill[gray!80] (\x,\y) circle (1.8pt);
    }
  }

  % =========== 原图形：长方形 (黑色) ===========
  % 宽为2，长(高)为4，面积 = 2 * 4 = 8
  % 左下角定于 (1, 1)
  \draw[black, thick] (1,1) rectangle (3,5);

  % =========== 解答图形1：平行四边形 (红色) ===========
  % 底为2，高为4，面积 = 2 * 4 = 8
  % 从 x=4 开始画底，底从 (4,1) 到 (6,1)
  % 上底整体向右平移 1 个单位，从 (5,5) 到 (7,5)
  \draw[red!70!black, line width=1.3pt] (4,1) -- (6,1) -- (7,5) -- (5,5) -- cycle;

  % =========== 解答图形2：三角形 (红色) ===========
  % 底为4，高为4，面积 = 4 * 4 / 2 = 8
  % 从 x=8 开始画底，底从 (8,1) 到 (12,1)
  % 顶点选取中间位置确保视觉对称，即 (10,5)
  \draw[red!70!black, line width=1.3pt] (8,1) -- (12,1) -- (10,5) -- cycle;

  % =========== 解答图形3：梯形 (红色) ===========
  % 上底1，下底3，高为4，面积 = (1 + 3) * 4 / 2 = 8
  % 从 x=14 开始画下底，下底从 (14,1) 到 (17,1)
  % 参考原答案图，左侧边垂直（直角梯形），所以上底从 (14,5) 开始
  % 上底长1，到 (15,5)
  \draw[red!70!black, line width=1.3pt] (14,1) -- (17,1) -- (15,5) -- (14,5) -- cycle;

\end{tikzpicture}
\end{center}

\vspace{0.8cm}

{\small\color{gray}
\textbf{解析：}\\
长方形的面积 = 长 $\times$ 宽 = $2 \times 4 = 8$。\\
题目要求画出的图形面积都需要为 $8$。为方便在点阵图作图，我们统一规定这三个图形的高都取 $4$（占用4个竖直单元格）。\\
1. \textbf{平行四边形的面积} = 底 $\times$ 高。已知面积=$8$，高=$4$，所以**底必须为 $2$**。\\
2. \textbf{三角形的面积} = 底 $\times$ 高 $\div\ 2$。已知面积=$8$，高=$4$，可列式：底 $\times 4 \div 2 = 8$，解得**底必须为 $4$**。\\
3. \textbf{梯形的面积} = (上底 + 下底) $\times$ 高 $\div\ 2$。已知面积=$8$，高=$4$，可列式：(上底 + 下底) $\times 4 \div 2 = 8$，解得 (上底 + 下底)= $4$。结合点阵特点，可以取**上底为 $1$，下底为 $3$**（直角梯形画法最清晰）。
}

\vspace{0.6cm}

{\small\color{blue!70!black}
\textbf{绘图基本原理与步骤：}\\
\textbf{1. 点子图阵列构建}：利用 TikZ 强大的 \texttt{\char`\\foreach} 嵌套循环功能，快速生成一整片均匀排列的圆点（\texttt{circle (1.8pt)}），作为辅助坐标系。\\
\textbf{2. 整数格点捕捉}：所有的图形顶点（如 \texttt{(4,1)}, \texttt{(7,5)}）都被精确设定在整数坐标上，确保线段端点能够地落在中国点子阵列图的几何视觉点上，严格遵循了“在点子图上作图”的要求。\\
\textbf{3. 色彩区分层级}：原图矩形选用黑色刻画题目条件，要求绘制的平行四边形、三角形、梯形统一选用视觉鲜明的深红色（\texttt{red!70!black}，线宽 \texttt{1.3pt}），突出作图结果。
}

\newpage


\section{解答：从上面看各是什么形状}

\vspace{0.4cm}

\noindent\textbf{1. 从上面看圆柱、正方体、圆锥和长方体，看到的各是什么形状？连一连。}

\vspace{1.2cm}

\begin{center}
% =====================================================================
% 【整体绘制思路：立体图形俯视图连线】
% 1. 坐标系与布局：
%    - 上排 (立体图形)：y=5 水平线。从左到右依次为 圆柱(x=2)、正方体(x=6)、圆锥(x=10)、长方体(x=15)。
%    - 下排 (平面图形)：y=1 水平线。从左到右依次为 长方形(x=2)、带圆心的圆(x=6)、圆(x=10)、正方形(x=15)。
% 2. 三维模拟绘画法：
%    - 圆柱：上下两个椭圆，两侧垂直连线。
%    - 正方体/长方体：绘制正面矩形，再通过固定偏移向量绘制顶面和侧面，形成等大透视感。
%    - 圆锥：底面画完整椭圆和虚线后半段，连接顶点。
% 3. 连线锚点与线型：
%    - 为每个图形的正下方(T)和正上方(B)预埋 \coordinate 锚点。
%    - 依据正确答案：圆柱连圆、正方体连正方形、圆锥连带点圆、长方体连长方形，使用红色粗线绘制。
% =====================================================================
\begin{tikzpicture}[>=Stealth, line width=1.0pt, scale=0.85]

  % =========== 上排：立体图形 ===========
  
  % 1. 圆柱 (Cylinder)
  \begin{scope}[shift={(2, 5)}]
    \draw (0, 0.8) ellipse (0.9 and 0.25);
    \draw (-0.9, 0.8) -- (-0.9, -0.6);
    \draw (0.9, 0.8) -- (0.9, -0.6);
    \draw (-0.9, -0.6) arc (180:360:0.9 and 0.25);
    % 预留连接锚点
    \coordinate (T1) at (0, -1.0);
  \end{scope}

  % 2. 正方体 (Cube)
  \begin{scope}[shift={(6, 5)}]
    % 正面 (尺寸 1.4 x 1.4)
    \draw (-0.7, -0.7) rectangle (0.7, 0.7);
    % 利用偏移量 (0.5, 0.5) 画顶面和右侧面
    \coordinate (FrontTL) at (-0.7, 0.7);
    \coordinate (FrontTR) at (0.7, 0.7);
    \coordinate (FrontBR) at (0.7, -0.7);
    \coordinate (BackTL)  at (-0.2, 1.2);
    \coordinate (BackTR)  at (1.2, 1.2);
    \coordinate (BackBR)  at (1.2, -0.2);
    
    \draw (FrontTL) -- (BackTL) -- (BackTR) -- (FrontTR);
    \draw (FrontTR) -- (BackTR) -- (BackBR) -- (FrontBR);
    \coordinate (T2) at (0, -1.0);
  \end{scope}

  % 3. 圆锥 (Cone)
  \begin{scope}[shift={(10, 5)}]
    % 底部椭圆弧
    \draw (-1.0, -0.6) arc (180:360:1.0 and 0.25);
    \draw[dashed] (1.0, -0.6) arc (0:180:1.0 and 0.25);
    % 连线到顶点
    \draw (-1.0, -0.6) -- (0, 1.3) -- (1.0, -0.6);
    \coordinate (T3) at (0, -1.0);
  \end{scope}

  % 4. 长方体 (Rectangular Cuboid)
  \begin{scope}[shift={(15, 5)}]
    % 正面 (尺寸 2.4 x 1.0，宽而扁)
    \draw (-1.2, -0.5) rectangle (1.2, 0.5);
    % 利用偏移量 (0.6, 0.5) 画顶面和右侧面
    \coordinate (FrontTL4) at (-1.2, 0.5);
    \coordinate (FrontTR4) at (1.2, 0.5);
    \coordinate (FrontBR4) at (1.2, -0.5);
    \coordinate (BackTL4)  at (-0.6, 1.0);
    \coordinate (BackTR4)  at (1.8, 1.0);
    \coordinate (BackBR4)  at (1.8, 0);
    
    \draw (FrontTL4) -- (BackTL4) -- (BackTR4) -- (FrontTR4);
    \draw (FrontTR4) -- (BackTR4) -- (BackBR4) -- (FrontBR4);
    \coordinate (T4) at (0.3, -1.0); % 因为向右延展，连线锚点稍微右移找中心
  \end{scope}


  % =========== 下排：平面视图 ===========

  % A. 长方形
  \begin{scope}[shift={(2, 1)}]
    \draw (-1.1, -0.6) rectangle (1.1, 0.6);
    \coordinate (B1) at (0, 1.2);
  \end{scope}

  % B. 带圆心的圆
  \begin{scope}[shift={(6, 1)}]
    \draw (0, 0) circle (0.9);
    \fill (0, 0) circle (1.5pt);
    \coordinate (B2) at (0, 1.2);
  \end{scope}

  % C. 普通圆
  \begin{scope}[shift={(10, 1)}]
    \draw (0, 0) circle (0.9);
    \coordinate (B3) at (0, 1.2);
  \end{scope}

  % D. 正方形
  \begin{scope}[shift={(15, 1)}]
    \draw (-0.7, -0.7) rectangle (0.7, 0.7);
    \coordinate (B4) at (0, 1.2);
  \end{scope}

  % =========== 绘制答案连线 (红色) ===========
  % \draw[color, line width] (起点) -- (终点);
  
  % 1. 圆柱从上面看是：普通圆 (C)
  \draw[red!80!black, line width=1.5pt] (T1) -- (B3);
  
  % 2. 正方体从上面看是：正方形 (D)
  \draw[red!80!black, line width=1.5pt] (T2) -- (B4);
  
  % 3. 圆锥从上面看是：带圆心的圆 (B)
  \draw[red!80!black, line width=1.5pt] (T3) -- (B2);
  
  % 4. 长方体从上面看是：长方形 (A)
  \draw[red!80!black, line width=1.5pt] (T4) -- (B1);

\end{tikzpicture}
\end{center}

\vspace{0.8cm}

{\small\color{gray}
\textbf{解析：}\\
观察立体图形的三视图（尤其是俯视图），想象从物体正上方往下看所能看到的轮廓平面图。\\
1. \textbf{圆柱}：上下底面都是圆，侧面垂直，所以从正上方往下看，只能看到一个与底面相同的**圆**。\\
2. \textbf{正方体}：所有的六个面都是完全相等的正方形，从正上方往下看，看到的就是其顶面，是一个**正方形**。\\
3. \textbf{圆锥}：底面是一个圆，最上面是一个尖顶。从正上方俯视时，能看到底面的圆形轮廓，并且能看到正中央那个尖顶（在平面图上体现为一个圆心点）。因此俯视图是**一个带有一点（圆心）的圆**。\\
4. \textbf{长方体}：相对的两个面是全等的长方形，从正上方往下看，看到的是它原本作为顶面的那个**长方形**。
}

\vspace{0.6cm}

{\small\color{blue!70!black}
\textbf{绘图基本原理与步骤：}\\
\textbf{1. 伪 3D 效果手工绘制}：纯利采用二维基础命令线构建极具透视感的各种立体图。比如正方形与长方体利用平行四边形偏移法画出顶面与侧边；圆柱和圆锥利用 \texttt{ellipse} (椭圆) 来呈现圆面在三维视角下的透视形变，圆锥隐蔽的背面底弧则贴心地使用了虚线 \texttt{dashed}。\\
\textbf{2. 预设解耦描点连线结构}：图形绘制采用相互独立、定距排列的 \texttt{scope} 分组模块化编写。并在每个模块的图形下方/上方埋下了不可见的透明锚点（如 \texttt{T1,T2} 以及 \texttt{B1,B2}）。当绘制连线时，仅需将对应的锚点对接即可，使得相交乱序的红线在代码层面上依然非常清爽清晰，且极易修正错题。\\
\textbf{3. 色彩重点剥离}：黑色的图作为客观存在的认知题干选项；红色加粗连线（\texttt{red!80!black}）强引力充当核心作答答案，对比鲜明符合辅导教材的美学要求。
}


\newpage
